
\subsection{Process trees}\label{sec:process_trees}

Process trees are a block-structured process-mining formalism widely used in discovery algorithms and model-quality evaluation~\cite{leemansDiscoveringBlockStructuredProcess2013,buijsQualityDimensionsProcess2014}. Each internal node carries one of four control-flow operators (sequence, XOR-choice, AND-parallel, redo) and the compositional structure guarantees soundness of the resulting workflow-net translation. We translate OC process trees (OCPTs) into decorated cospans by reading interface sets directly from the operator structure, following the same cospan-algebra construction as for OCPNs and OC causal nets.

\begin{definition}[Object-centric process tree]
  \label{def:process-tree}
  An object-centric (OC) process tree~\cite{leemansDiscoveringBlockStructuredProcess2013,vandettenDiscoveringCompactLive2024}
  over a finite set of object types $\mathcal{O}$ is a finite rooted tree of the following form.
  \begin{itemize}
    \item Each internal node is a control-flow operator, applied to its child subtrees, one of
      sequence $\to$, exclusive choice $\times$, concurrency $+$, or redo $\circlearrowleft$.
    \item Each leaf is an object-centric activity $(a,\mathrm{rel})$ with an activity name $a$ and
      the object types $\mathrm{rel}\subseteq\mathcal{O}$ it involves.
  \end{itemize}
  The open ends $\bot$ and $\top$ of Definition~\ref{def:contexts-general} bound the root, and the
  analysis may attach a start generator $\gamma_1$ and an end generator $\gamma_2$ there
  (Section~\ref{subsec:analysis}). Each object type runs
  through the operators on its own, independently of the others.
\end{definition}

\begin{figure}[tb]
  \centering
  \newcommand{\pwire}[2]{%
  \draw[bluewire] let \p1=($(#2.center)-(#1.center)$), \n1={atan2(\y1,\x1)} in
    ($(#1.\n1)!1.9pt!90:(#2.center)$) -- ($(#2.\n1+180)!1.9pt!-90:(#1.center)$);
  \draw[redwire]  let \p1=($(#2.center)-(#1.center)$), \n1={atan2(\y1,\x1)} in
    ($(#1.\n1)!1.9pt!-90:(#2.center)$) -- ($(#2.\n1+180)!1.9pt!90:(#1.center)$);}
\begin{tikzpicture}[
    op/.style={circle, draw=black!55, fill=black!4, inner sep=1.2pt, minimum size=5mm,
               font=\small},
    lboth/.style={rectangle, rounded corners=2pt, draw=black!55, fill=black!5,
                  minimum size=6mm, font=\small},
    lord/.style={rectangle, rounded corners=2pt, draw=bluecol, fill=bluecol!12,
                 minimum size=6mm, font=\small},
    litem/.style={rectangle, rounded corners=2pt, draw=redcol, fill=redcol!12,
                  minimum size=6mm, font=\small}]
  \node[op]    (root) at (0,1.8)     {$\to$};
  \node[lboth] (a)    at (-1.6,0.6)  {$a$};
  \node[op]    (plus) at (0,0.6)     {$+$};
  \node[lboth] (s)    at (1.6,0.6)   {$s$};
  \node[lord]  (b)    at (-0.65,-0.6){$b$};
  \node[litem] (c)    at (0.65,-0.6) {$c$};
  \pwire{root}{a}
  \pwire{root}{plus}
  \pwire{root}{s}
  \draw[bluewire] (plus) -- (b);
  \draw[redwire]  (plus) -- (c);
\end{tikzpicture}
\unskip
  \caption{The OCPT example $\to(a,+(b,c),s)$, with tree connectors coloured by object type (order
    blue, item red). A connector carrying both types is two parallel coloured wires, a single type
    one coloured wire.}
  \label{fig:re-ocpt}
\end{figure}

The operators apply per object. $\to(S_1,\ldots,S_n)$ runs its subtrees in order.
$\times(S_1,\ldots,S_n)$ runs exactly one of them. Under $+(S_1,\ldots,S_n)$ they all run
concurrently. Finally, $\circlearrowleft(S,B_1,\ldots,B_m)$ runs $S$ once, then any number of times
takes some redo $B_j$ and runs $S$ again. Each object type follows its own path, so one type may
take a $\times$-branch another does not, and a type outside a leaf's $\mathrm{rel}$ skips that
leaf. Leaf labels may repeat, each occurrence firing independently. We write $T$ for the leaf
activities and $T_\tau\subseteq T$ for the silent leaves. A silent leaf carries a type but no
observable label and models a skip, the XOR alternative $\times(a,\tau)$ to the activity it
bypasses. The induced typed directly-follows arcs are $D_{PT}\subseteq T\times\mathcal{O}\times T$,
with $(x,\omega,t)\in D_{PT}$ when type $\omega$ flows from leaf $x$ to leaf $t$. When a leaf
relates to a single type $\mathrm{type}(t)$, the arcs $D_{PT}$ are forced by the tree. The source
formalism also marks a leaf's types as divergent, convergent or deficient. Those marks are
existential statements over an event log~\cite{vandettenDiscoveringCompactLive2024}, properties of
the flattened object-centric data a tree is discovered from, and the cospan language needs none of
them. A count constraint on a leaf's legs is written in $\Lambda$, as for every notation
(Definition~\ref{def:multiplicity-decoration}).

This construction instantiates the general framework of Section~\ref{sec:general-framework}, and
the two mapping choices are simple. An observable leaf is an activity. Each control-flow operator
becomes a mediator, the choice $\times$ an XOR and the parallel $+$ an AND. Sequence $\to$ and redo
$\circlearrowleft$ need no mediator of their own. They are built from these two, and a
degree-$(1,1)$ mediator, where AND, XOR and a SEQ pass-through all coincide, simply forwards its one
branch. A silent leaf is a transparent mediator and contributes no generator, as for the silent
transitions of Section~\ref{sec:petri-nets}. Compiling the operator tree this way gives an AND/OR
graph, the compiled flow graph $\mathcal{G}_{PT}$ (Definition~\ref{def:pt-flow-graph}), over which
the construction is the walk of Definition~\ref{def:contexts-general}. This is the instance where
the walk is not immediate. Where the object-centric causal net pre-compiles its routing into binding
sets and the walk ends at once, the process tree spreads its routing across operator structure that
the walk must cross. The instance map is as follows.
\medskip
\begin{center}
  \small\begin{tabularx}{\linewidth}{@{}lX@{}}
    \hline
    General framework & Process tree \\
    \hline
    AND/OR graph $(V,E)$ & the compiled flow graph $\mathcal{G}_{PT}$ (Definition~\ref{def:pt-flow-graph}) \\
    Activities $\mathcal{A}$ & observable leaves $T\setminus T_\tau$ \\
    AND mediators & the parallel operator $+$ (split and join) \\
    XOR mediators & the choice operator $\times$, and a loop's redo-or-exit \\
    SEQ pass-throughs & the sequence operator $\to$ and the silent leaves $T_\tau$ (transparent); the degree-$(1,1)$ case where AND, XOR and SEQ coincide \\
    \hline
  \end{tabularx}
\end{center}
\medskip
\begin{definition}[Compiled flow graph]
  \label{def:pt-flow-graph}
  The compiled flow graph $\mathcal{G}_{PT}=\mathcal{G}(\Omega)$, for $\Omega$ the operator tree of $PT$, is built bottom-up from the
  operator tree, the map $\mathcal{G}(\cdot)$ sending each subtree to a directed graph with one
  entry and one exit. It is the AND/OR graph $(V,E)$ of the instance map, not the raw
  parent-to-child tree. The cases are as follows.
  \begin{itemize}
    \item An observable leaf $t\in T\setminus T_\tau$ gives the single activity node $t$.
    \item A silent leaf $\tau\in T_\tau$ gives a degree-$(1,1)$ SEQ pass-through carrying type
      $\mathrm{type}(\tau)$, transparent to the walk.
    \item For $\to(S_1,\ldots,S_n)$, chain the children through $n{-}1$ fresh SEQ pass-throughs,
      $\mathcal{G}(S_1)\to s_1\to\cdots\to s_{n-1}\to\mathcal{G}(S_n)$. The ordering is carried by the
      chain, not by any one node.
    \item For $+(S_1,\ldots,S_n)$, a fresh AND-split $p^+$ fans into each $\mathcal{G}(S_i)$ and a
      fresh AND-join $q^+$ collects them.
    \item For $\times(S_1,\ldots,S_n)$, take a fresh XOR-split $x^\times$ and XOR-join $y^\times$, each
      branch $\mathcal{G}(S_i)$ running between them. The walk keeps the reached sets of the branches separate at $x^\times$.
    \item For $\circlearrowleft(S,B_1,\ldots,B_m)$, take a mandatory $\mathcal{G}(S)$, then an XOR mediator
      that either exits to the next block or takes a back-edge through some $\mathcal{G}(B_j)$ and
      repeats $S$. A repeated pass of the loop is a composite of generators and adds no signature
      element.
  \end{itemize}
\end{definition}

\begin{figure*}[tp]
  \centering
  \begin{tikzpicture}[petriBase, scale=0.95, inline,
    task/.style={rectangle, rounded corners=3pt, draw=black!60, fill=black!4,
                 minimum size=7mm, font=\small},
    btask/.style={task, draw=bluecol, fill=bluecol!6},
    rtask/.style={task, draw=redcol,  fill=redcol!6},
    agw/.style={diamond, draw=black!60, fill=black!4, inner sep=1pt, minimum size=6mm,
                font=\footnotesize}]
  \node[task]  (a)  at (0,0)      {$a$};
  \node[agw]   (ps) at (1.9,0)    {$+$};
  \node[btask] (b)  at (3.9,0.85) {$b$};
  \node[rtask] (c)  at (3.9,-0.85){$c$};
  \node[agw]   (qj) at (5.9,0)    {$+$};
  \node[task]  (s)  at (7.8,0)    {$s$};
  \draw[blueflow] ([yshift=1.4mm]a.east)  -- (1.74,0.14);
  \draw[redflow]  ([yshift=-1.4mm]a.east) -- (1.74,-0.14);
  \draw[blueflow] (ps) -- (b);
  \draw[redflow]  (ps) -- (c);
  \draw[blueflow] (b)  -- (qj);
  \draw[redflow]  (c)  -- (qj);
  \draw[blueflow] (6.06,0.14)  -- ([yshift=1.4mm]s.west);
  \draw[redflow]  (6.06,-0.14) -- ([yshift=-1.4mm]s.west);
\end{tikzpicture}
\unskip
  \caption{The compiled flow graph $\mathcal{G}_{PT}$ of the example tree $\to(a,+(b,c),s)$
    (Fig.~\ref{fig:re-ocpt}). The $+$ node becomes an AND-split and AND-join ($+$) around the
    two concurrent typed lifecycles, $b$ on the order and $c$ on the item. The sequence links are
    degree-$(1,1)$ pass-throughs and collapse onto the edges.}
  \label{fig:re-pt-flow}
\end{figure*}

A silent leaf is needed exactly where a typed token crosses a span of the operator tree with
no observable activity on it. Where an observable leaf already lies on that span, none is
needed. The engine of Section~\ref{sec:general-framework} then runs unchanged, with one
instance-specific fact. The walk traverses the compiled flow graph $\mathcal{G}_{PT}$ rather
than the raw tree, keeping the branches of each XOR-split separate, crossing silent leaves
transparently, and stopping at the first activity on each side of a $\circlearrowleft$ back-edge,
so no unrolling is needed. Branches that leave the root end at the open ends $\bot$ and $\top$, and
start and end generators there are a choice of the analysis (Section~\ref{subsec:analysis}). Each context $(P,S)$ of a leaf $t$
yields a generator on the typed wires of $D_{PT}$. Each leg carries a count, by default pinned to
$n_p=1$ (one object per firing), and any other constraint on the leaf's own legs enters its
constraint system $\Lambda$ as a per-leg interval or a linear equality or inequality among its counts.
Generators sharing a typed wire compose along it in $\mathbf{F}(\Sigma)$ (Definition~\ref{def:free-category}).

\begin{theorem}[Canonical cospan-algebra presentation of OC process trees]
  \label{thm:ptree-to-cospan-signature}
  Let $PT$ be an OC process tree (Definition~\ref{def:process-tree}). Then $PT$ canonically determines a cospan-algebra signature $\Sigma$ in which each typed wire $(x,w,t)$ carries the set $w$ of object types $\omega$ with $(x,\omega,t)\in D_{PT}$ on its path, and hence the free symmetric monoidal category $\mathbf{F}(\Sigma)$.
\end{theorem}

\begin{proof}
  The schema of Section~\ref{sec:notation-by-notation} applies at the instance map above,
  finiteness of $T$ and of the operator tree $\Omega$ giving the conditions of Definition~\ref{def:lm-graph}, with
  each edge typed by the arc types in~$D_{PT}$. The delta concerns the loops. The walk runs on the
  compiled flow graph $\mathcal{G}_{PT}$ (Definition~\ref{def:pt-flow-graph}), which is finite, and a
  branch stops at any node already on its path, so every reach family and every set of contexts is
  finite (Definition~\ref{def:contexts-general}). A further pass of a $\circlearrowleft$ is a
  composite of generators already present and adds none, and the firing rule is self-dual across
  all four operators.
\end{proof}

\begin{figure*}[tp]
  \centering
  \begin{tikzpicture}[inline,
    bspd/.style={spider, fill=bluecol},
    rspd/.style={spider, fill=redcol},
    mbox/.style={draw=black, rounded corners=4pt, fill=white, minimum size=0.75cm, font=\small},
    plabel/.style={font=\tiny}]

  \node[font=\footnotesize] at (2.7,1.4) {decorated cospan};
  \node[font=\footnotesize] at (9.0,1.4) {string diagram};

  \node[font=\small] at (-2.3,0) {$g_a$};
  \node[cospanblue, minimum height=0.9cm] (aL) at (0,0) {};
  \node[bspd] at (0,0.20) {};
  \node[rspd] at (0,-0.20) {};
  \node at (0.9,0) {$\rightarrow$};
  \node[cospangrey, minimum height=1.0cm, minimum width=2.6cm] (aA) at (2.7,0) {};
  \node[mbox, minimum height=0.8cm] (af) at (2.7,0) {$a$};
  \node[bspd] (aao) at (1.65,0.20) {};
  \node[rspd] (aai) at (1.65,-0.20) {};
  \node[bspd] (abo) at (3.75,0.20) {};
  \node[rspd] (abi) at (3.75,-0.20) {};
  \draw[bluewire] (aao) -- (af.160);
  \draw[redwire]  (aai) -- (af.200);
  \draw[bluewire] (af.20)  -- (abo);
  \draw[redwire]  (af.-20) -- (abi);
  \node at (4.5,0) {$\leftarrow$};
  \node[cospanblue, minimum height=0.9cm] (aR) at (5.4,0) {};
  \node[bspd] at (5.4,0.20) {};
  \node[rspd] at (5.4,-0.20) {};
  \node[plabel, anchor=east] at ([shift={(-3pt,0.20cm)}]aL.west) {$(\bot,\mathrm{o},a)$};
  \node[plabel, anchor=east] at ([shift={(-3pt,-0.20cm)}]aL.west) {$(\bot,\mathrm{it},a)$};
  \node[plabel, anchor=west] at ([shift={(3pt,0.20cm)}]aR.east) {$(a,\mathrm{o},b)$};
  \node[plabel, anchor=west] at ([shift={(3pt,-0.20cm)}]aR.east) {$(a,\mathrm{it},c)$};
  \node at (7.4,0) {$\mapsto$};
  \node[morphismblue, minimum height=0.95cm] (asd) at (9.0,0) {$a$};
  \draw[bluewire] ($(asd.south west)!0.70!(asd.north west)$) -- ++(-0.55,0);
  \draw[redwire]  ($(asd.south west)!0.30!(asd.north west)$) -- ++(-0.55,0);
  \draw[bluewire] ($(asd.south east)!0.70!(asd.north east)$) -- ++(0.55,0);
  \draw[redwire]  ($(asd.south east)!0.30!(asd.north east)$) -- ++(0.55,0);

  \begin{scope}[yshift=-1.5cm]
    \node[font=\small] at (-2.3,0) {$g_b$};
    \node[cospanblue, minimum height=0.6cm] (bL) at (0,0) {};
    \node[bspd] at (0,0) {};
    \node at (0.9,0) {$\rightarrow$};
    \node[cospangrey, minimum height=1.0cm, minimum width=2.6cm] (bA) at (2.7,0) {};
    \node[mbox] (bf) at (2.7,0) {$b$};
    \node[bspd] (bao) at (1.65,0) {};
    \node[bspd] (bbo) at (3.75,0) {};
    \draw[bluewire] (bao) -- (bf.west);
    \draw[bluewire] (bf.east) -- (bbo);
    \node at (4.5,0) {$\leftarrow$};
    \node[cospanblue, minimum height=0.6cm] (bR) at (5.4,0) {};
    \node[bspd] at (5.4,0) {};
    \node[plabel, anchor=east] at ([shift={(-3pt,0cm)}]bL.west) {$(a,\mathrm{o},b)$};
    \node[plabel, anchor=west] at ([shift={(3pt,0cm)}]bR.east) {$(b,\mathrm{o},s)$};
    \node at (7.4,0) {$\mapsto$};
    \node[morphismblue, minimum size=0.75cm] (bsd) at (9.0,0) {$b$};
    \draw[bluewire] (bsd.west) -- ++(-0.55,0);
    \draw[bluewire] (bsd.east) -- ++(0.55,0);
  \end{scope}

  \begin{scope}[yshift=-3.0cm]
    \node[font=\small] at (-2.3,0) {$g_c$};
    \node[cospanblue, minimum height=0.6cm] (cL) at (0,0) {};
    \node[rspd] at (0,0) {};
    \node at (0.9,0) {$\rightarrow$};
    \node[cospangrey, minimum height=1.0cm, minimum width=2.6cm] (cA) at (2.7,0) {};
    \node[mbox] (cf) at (2.7,0) {$c$};
    \node[rspd] (cai) at (1.65,0) {};
    \node[rspd] (cbi) at (3.75,0) {};
    \draw[redwire] (cai) -- (cf.west);
    \draw[redwire] (cf.east) -- (cbi);
    \node at (4.5,0) {$\leftarrow$};
    \node[cospanblue, minimum height=0.6cm] (cR) at (5.4,0) {};
    \node[rspd] at (5.4,0) {};
    \node[plabel, anchor=east] at ([shift={(-3pt,0cm)}]cL.west) {$(a,\mathrm{it},c)$};
    \node[plabel, anchor=west] at ([shift={(3pt,0cm)}]cR.east) {$(c,\mathrm{it},s)$};
    \node at (7.4,0) {$\mapsto$};
    \node[morphismblue, minimum size=0.75cm] (csd) at (9.0,0) {$c$};
    \draw[redwire] (csd.west) -- ++(-0.55,0);
    \draw[redwire] (csd.east) -- ++(0.55,0);
  \end{scope}

  \begin{scope}[yshift=-4.5cm]
    \node[font=\small] at (-2.3,0) {$g_s$};
    \node[cospanblue, minimum height=0.9cm] (sL) at (0,0) {};
    \node[bspd] at (0,0.20) {};
    \node[rspd] at (0,-0.20) {};
    \node at (0.9,0) {$\rightarrow$};
    \node[cospangrey, minimum height=1.0cm, minimum width=2.6cm] (sA) at (2.7,0) {};
    \node[mbox, minimum height=0.8cm] (sf) at (2.7,0) {$s$};
    \node[bspd] (sao) at (1.65,0.20) {};
    \node[rspd] (sai) at (1.65,-0.20) {};
    \node[bspd] (sbo) at (3.75,0.20) {};
    \node[rspd] (sbi) at (3.75,-0.20) {};
    \draw[bluewire] (sao) -- (sf.160);
    \draw[redwire]  (sai) -- (sf.200);
    \draw[bluewire] (sf.20)  -- (sbo);
    \draw[redwire]  (sf.-20) -- (sbi);
    \node at (4.5,0) {$\leftarrow$};
    \node[cospanblue, minimum height=0.9cm] (sR) at (5.4,0) {};
    \node[bspd] at (5.4,0.20) {};
    \node[rspd] at (5.4,-0.20) {};
    \node[plabel, anchor=east] at ([shift={(-3pt,0.20cm)}]sL.west) {$(b,\mathrm{o},s)$};
    \node[plabel, anchor=east] at ([shift={(-3pt,-0.20cm)}]sL.west) {$(c,\mathrm{it},s)$};
    \node[plabel, anchor=west] at ([shift={(3pt,0.20cm)}]sR.east) {$(s,\mathrm{o},\top)$};
    \node[plabel, anchor=west] at ([shift={(3pt,-0.20cm)}]sR.east) {$(s,\mathrm{it},\top)$};
    \node at (7.4,0) {$\mapsto$};
    \node[morphismblue, minimum height=0.95cm] (ssd) at (9.0,0) {$s$};
    \draw[bluewire] ($(ssd.south west)!0.70!(ssd.north west)$) -- ++(-0.55,0);
    \draw[redwire]  ($(ssd.south west)!0.30!(ssd.north west)$) -- ++(-0.55,0);
    \draw[bluewire] ($(ssd.south east)!0.70!(ssd.north east)$) -- ++(0.55,0);
    \draw[redwire]  ($(ssd.south east)!0.30!(ssd.north east)$) -- ++(0.55,0);
  \end{scope}
\end{tikzpicture}
\unskip
  \caption{The generator cospans of Fig.~\ref{fig:re-ocpt}, each as a decorated cospan and the
    string diagram it denotes. $g_a$ takes the order and item threads from the open end $\bot$, $g_b,g_c$ carry them, and
    $g_s$ synchronises both.}
  \label{fig:re-ocpt-cospans}
\end{figure*}

\begin{exmp}[A small OCPT and its generator cospans]
  \label{ex:running-ptree-cospan}
  Consider the OC process tree $\to(a, +(b,c), s)$ of Fig.~\ref{fig:re-ocpt}, compiled to the
  flow graph of Fig.~\ref{fig:re-pt-flow}, over the object types $\mathrm{order}$ and
  $\mathrm{item}$, with typed dependency arcs
  \[
    D_{PT} = \{(a,\mathrm{o},b),\,(a,\mathrm{it},c),\,
      (b,\mathrm{o},s),\,(c,\mathrm{it},s)\}.
  \]
  The $+$ node carries the object-centricity. It runs two concurrent typed
  object lifecycles, the order through $b$ and the item through $c$. Reading a
  generator off each leaf context (Fig.~\ref{fig:re-ocpt-cospans}) gives the first activity
  $g_a:\{(\bot,\mathrm{o},a),(\bot,\mathrm{it},a)\}\to\{(a,\mathrm{o},b),(a,\mathrm{it},c)\}$, the two lifecycle
  threads $g_b:\{(a,\mathrm{o},b)\}\to\{(b,\mathrm{o},s)\}$ and
  $g_c:\{(a,\mathrm{it},c)\}\to\{(c,\mathrm{it},s)\}$, and the synchronisation
  \begin{gather*}
    g_s = \Bigl(\{(b,\mathrm{o},s),(c,\mathrm{it},s)\}
      \xrightarrow{\iota} A \xleftarrow{\iota}\twocolbreak  \{(s,\mathrm{o},\top),(s,\mathrm{it},\top)\},\; d_s\Bigr),
  \end{gather*}
  whose left boundary carries both object types. Each typed wire here carries a single type,
  the order thread staying order-typed and the item thread item-typed, with $a$ and $s$ the only
  generators touching both. The two lifecycles are independent until $s$, which is exactly the
  concurrency the $+$ node expresses. Every leg here carries the default count of one.
  This concurrency is exactly what the certificate separates from exclusive choice (Lemma~\ref{lem:sig-complete}).
\end{exmp}

\begin{figure}[tb]
  \centering
  \newcommand{\pwire}[3]{%
  \draw[bluewire,#1] let \p1=($(#3.center)-(#2.center)$), \n1={atan2(\y1,\x1)} in
    ($(#2.\n1)!1.9pt!90:(#3.center)$) -- ($(#3.\n1+180)!1.9pt!-90:(#2.center)$);
  \draw[redwire,#1]  let \p1=($(#3.center)-(#2.center)$), \n1={atan2(\y1,\x1)} in
    ($(#2.\n1)!1.9pt!-90:(#3.center)$) -- ($(#3.\n1+180)!1.9pt!90:(#2.center)$);}
\begin{tikzpicture}[
    op/.style={circle, draw=black!55, fill=black!4, inner sep=1.2pt, minimum size=5.5mm,
               font=\small},
    leaf/.style={rectangle, rounded corners=2pt, draw=black!55, fill=black!5,
                 minimum size=6mm, font=\small},
    elab/.style={font=\tiny, black!60}]
  \node[op]   (root) at (0,1.5)     {$\circlearrowleft$};
  \node[leaf] (a)    at (-1.7,0)    {$a$};
  \node[leaf] (b)    at (0,0)       {$b$};
  \node[leaf] (c)    at (1.7,0)     {$c$};
  \pwire{solid}{root}{a}
  \pwire{dashed}{root}{b}
  \pwire{dashed}{root}{c}
  \node[elab] at (-1.15,0.85) {do};
  \node[elab, anchor=west] at (0.14,0.75) {redo};
  \node[elab] at (1.42,0.85)  {redo};
  \node[font=\tiny, black!60, align=center] at (0,-0.75)
    {every leaf: order $n=1$, item $n=1$};
\end{tikzpicture}
\unskip
  \caption{The redo tree $\circlearrowleft(a,b,c)$, a mandatory body $a$ (solid ``do'' edge) and two
    redo bodies $b,c$ (dashed ``redo'' edges). Every leaf relates to both object types, order (blue) and
    item (red), one of each per firing, so each connector is two parallel coloured wires.}
  \label{fig:re-loop-tree}
\end{figure}

\begin{figure}[tb]
  \centering
  \begin{tikzpicture}[petriBase, scale=1.0, inline,
    evt/.style={circle, draw=black!70, fill=black!6, minimum size=6.5mm, font=\small},
    task/.style={rectangle, rounded corners=3pt, draw=black!60, fill=black!4,
                 minimum size=7mm, font=\small},
    xgw/.style={diamond, draw=black!60, fill=black!4, inner sep=1pt, minimum size=6.5mm, font=\footnotesize}]
  \node[evt]  (g1) at (0,0)      {$\bot$};
  \node[task] (a)  at (2.0,0)    {$a$};
  \node[xgw]  (x)  at (4.0,0)    {$\times$};
  \node[evt]  (g2) at (6.1,0)    {$\top$};
  \node[task] (b)  at (3.0,-1.5) {$b$};
  \node[task] (c)  at (3.0,-3.1) {$c$};

  \draw[blueflow] ([yshift=2.3pt]g1.east) -- ([yshift=2.3pt]a.west);
  \draw[redflow]  ([yshift=-2.3pt]g1.east) -- ([yshift=-2.3pt]a.west);
  \draw[blueflow] ([yshift=2.3pt]a.east)  -- (3.76,0.081);
  \draw[redflow]  ([yshift=-2.3pt]a.east) -- (3.76,-0.081);
  \draw[blueflow] (4.24,0.081)  -- ([yshift=2.3pt]g2.west) node[pos=0.55, above=1pt, font=\tiny] {exit};
  \draw[redflow]  (4.24,-0.081) -- ([yshift=-2.3pt]g2.west);

  \draw[blueflow] (x.235) to[out=235,in=75]  ([xshift=-2.3pt]b.north);
  \draw[redflow]  (x.265) to[out=235,in=75]  ([xshift=2.3pt]b.north);
  \draw[blueflow] ([yshift=2.3pt]b.west)  to[out=155,in=-80] ([xshift=2.3pt]a.south);
  \draw[redflow]  ([yshift=-2.3pt]b.west) to[out=155,in=-80] ([xshift=-2.3pt]a.south);

  \draw[blueflow] (x.290) to[out=290,in=25]  ([yshift=2.3pt]c.east);
  \draw[redflow]  (x.320) to[out=290,in=25]  ([yshift=-2.3pt]c.east);
  \draw[blueflow] ([yshift=2.3pt]c.west)  to[out=170,in=-118] ([xshift=-3pt]a.south);
  \draw[redflow]  ([yshift=-2.3pt]c.west) to[out=170,in=-118] ([xshift=-7pt]a.south);

  \node[font=\tiny, align=center, above=1pt of a] {order $n=1$\\ item $n=1$};
  \node[font=\tiny, black!60, below=2pt of b] {redo body};
  \node[font=\tiny, black!60, below=2pt of c] {redo body};
\end{tikzpicture}
\unskip
  \caption{The compiled flow graph of $\circlearrowleft(a,b,c)$. After $a$, one XOR redo mediator
    either exits to the block after the loop or takes one of two back-edges, through redo body $b$ or
    through redo body $c$, and repeats $a$. Each connector carries both types as parallel coloured
    wires.}
  \label{fig:re-loop}
\end{figure}

\begin{figure*}[tp]
  \centering
  \newcommand{\dwire}[2]{%
  \draw[blueflow] ([yshift=2.1pt]#1.east) -- ([yshift=2.1pt]#2.west);%
  \draw[redflow]  ([yshift=-2.1pt]#1.east) -- ([yshift=-2.1pt]#2.west);}
\begin{tikzpicture}[petriBase, scale=0.85, inline,
    evt/.style={circle, draw=black!70, fill=black!6, minimum size=6mm, font=\small},
    task/.style={rectangle, rounded corners=3pt, draw=black!60, fill=black!4,
                 minimum size=6.5mm, font=\small, inner sep=2.5pt},
    dlab/.style={font=\footnotesize, anchor=east, black!70}]

  \node[font=\tiny, black!70, align=left, anchor=west] at (-0.9,1.05)
    {every pass: order $n=1$, item $n=1$};

  \node[dlab] at (-0.9,0) {depth $0$};
  \node[evt]  (a0) at (0,0)   {$\bot$};
  \node[task] (b0) at (1.6,0) {$a$};
  \node[evt]  (e0) at (3.2,0) {$\top$};
  \dwire{a0}{b0} \dwire{b0}{e0}

  \node[dlab] at (-0.9,-1.5) {depth $1$};
  \node[evt]  (a1) at (0,-1.5)   {$\bot$};
  \node[task] (b1) at (1.6,-1.5) {$a$};
  \node[task] (r1) at (3.2,-1.5) {$b$};
  \node[task] (c1) at (4.8,-1.5) {$a$};
  \node[evt]  (e1) at (6.4,-1.5) {$\top$};
  \dwire{a1}{b1} \dwire{b1}{r1} \dwire{r1}{c1} \dwire{c1}{e1}

  \node[dlab] at (-0.9,-3.0) {depth $2$};
  \node[evt]  (a2) at (0,-3.0)    {$\bot$};
  \node[task] (b2) at (1.6,-3.0)  {$a$};
  \node[task] (r2) at (3.2,-3.0)  {$b$};
  \node[task] (c2) at (4.8,-3.0)  {$a$};
  \node[task] (s2) at (6.4,-3.0)  {$c$};
  \node[task] (d2) at (8.0,-3.0)  {$a$};
  \node[evt]  (e2) at (9.6,-3.0)  {$\top$};
  \dwire{a2}{b2} \dwire{b2}{r2} \dwire{r2}{c2} \dwire{c2}{s2} \dwire{s2}{d2} \dwire{d2}{e2}
  \draw[decorate, decoration={brace, amplitude=4pt, mirror}, black!55]
    ($(r2.south west)+(0,-0.12)$) -- ($(s2.south east)+(0,-0.12)$)
    node[midway, below=4pt, font=\scriptsize, black!70]
    {$a$'s context between two redo bodies};
\end{tikzpicture}
\unskip
  \caption{The redo $\circlearrowleft(a,b,c)$ unrolled to depths $0$, $1$, and $2$, both types
    threaded per pass, one order and one item. The depth-$2$ unroll takes
    both redo bodies once, $a\to b\to a\to c\to a$. The middle $a$ then sits between two distinct
    redo bodies, one of its contexts, and deeper unrolling composes only generators already present.}
  \label{fig:re-pt-unroll}
\end{figure*}

\begin{exmp}[A redo, unrolled to its cospans]
  \label{ex:running-ptree-loop}
  The examples above are loop-free. This one shows the construction on a cycle. Take the redo tree $\circlearrowleft(a,b,c)$ over two object types,
  $\mathrm{order}$ (blue) and $\mathrm{item}$ (red) (Fig.~\ref{fig:re-loop-tree}). The mandatory body
  $a$ runs once, then the redo either exits or takes one of the two redo bodies $b,c$ and runs $a$
  again, any number of times. Every leaf relates to both types, one order and one item per firing, the default
  multiplicity. The compiled flow graph
  (Definition~\ref{def:pt-flow-graph}, Fig.~\ref{fig:re-loop}) joins $a,b,c$ by one XOR redo mediator
  with two back-edges, through $b$ and through $c$.

  Unrolling the redo makes the finite generator set visible (Fig.~\ref{fig:re-pt-unroll}). The
  depth-$2$ unroll takes both redo bodies once, $a\to b\to a\to c\to a$. The middle $a$ then sits
  between two distinct redo bodies, a redo body on the input side and one on the output for each
  type, which is one of its contexts. Deeper unrolling only composes generators already present, so
  the loop reduces to finitely many generators, $g_b$, $g_c$ and one $g_a$ per pairing of a predecessor ($\bot$, $b$ or $c$) with a successor ($b$, $c$ or $\top$), which the walk from $a$ finds by
  stopping at the first activity past the redo mediator on each side.

  Reading a generator off each leaf context (Fig.~\ref{fig:re-loop-cospans}) gives the loop-entry
  $g_a$ and the two redo-body generators $g_b,g_c$, each threading one order and
  one item. Repeating the loop composes $g_a$ with the redo blocks
  $g_b\,g_a$ and $g_c\,g_a$ in any order and adds no signature element. The loop language is
  recovered at the composition level, and the typed redo needs no new machinery.
\end{exmp}

\begin{figure*}[tp]
  \centering
  \resizebox{\textwidth}{!}{\begin{tikzpicture}[inline,
    bspd/.style={spider, fill=bluecol},
    rspd/.style={spider, fill=redcol},
    mbox/.style={draw=black, rounded corners=4pt, fill=white, minimum size=0.75cm, font=\small},
    plabel/.style={font=\tiny}]

  \node[font=\footnotesize] at (2.7,1.7) {decorated cospan};
  \node[font=\footnotesize] at (9.0,1.7) {string diagram};

  \foreach \row/\g/\src/\tgt in {0/a/b/b, -1.6/b/a/a, -3.2/c/a/a}{%
    \begin{scope}[yshift=\row cm]
      \node[font=\small] at (-2.3,0) {$g_{\g}$};
      \node[cospanblue, minimum height=0.9cm] (L) at (0,0) {};
      \node[bspd] at (0,0.20) {};
      \node[rspd] at (0,-0.20) {};
      \node at (0.9,0) {$\rightarrow$};
      \node[cospangrey, minimum height=1.0cm, minimum width=2.6cm] (A) at (2.7,0) {};
      \node[mbox, minimum height=0.8cm] (f) at (2.7,0) {$\g$};
      \node[bspd] (ao) at (1.65,0.20) {};
      \node[rspd] (ai) at (1.65,-0.20) {};
      \node[bspd] (bo) at (3.75,0.20) {};
      \node[rspd] (bi) at (3.75,-0.20) {};
      \draw[bluewire] (ao) -- (f.160);
      \draw[redwire]  (ai) -- (f.200);
      \draw[bluewire] (f.20)  -- (bo);
      \draw[redwire]  (f.-20) -- (bi);
      \node at (4.5,0) {$\leftarrow$};
      \node[cospanblue, minimum height=0.9cm] (R) at (5.4,0) {};
      \node[bspd] at (5.4,0.20) {};
      \node[rspd] at (5.4,-0.20) {};
      \node[plabel, anchor=east] at ([shift={(-3pt,0.20cm)}]L.west) {$(\src,\mathrm{o},\g)$};
      \node[plabel, anchor=east] at ([shift={(-3pt,-0.20cm)}]L.west) {$(\src,\mathrm{it},\g)$};
      \node[plabel, anchor=west] at ([shift={(3pt,0.20cm)}]R.east) {$(\g,\mathrm{o},\tgt)$};
      \node[plabel, anchor=west] at ([shift={(3pt,-0.20cm)}]R.east) {$(\g,\mathrm{it},\tgt)$};
      \node at (7.4,0) {$\mapsto$};
      \node[morphismblue, minimum height=0.95cm] (sd) at (9.0,0) {$\g$};
      \draw[bluewire] ($(sd.south west)!0.70!(sd.north west)$) -- ++(-0.55,0);
      \draw[redwire]  ($(sd.south west)!0.30!(sd.north west)$) -- ++(-0.55,0);
      \draw[bluewire] ($(sd.south east)!0.70!(sd.north east)$) -- ++(0.55,0);
      \draw[redwire]  ($(sd.south east)!0.30!(sd.north east)$) -- ++(0.55,0);
    \end{scope}}
\end{tikzpicture}
\unskip}
  \caption{The generator cospans of $\circlearrowleft(a,b,c)$, each as a decorated cospan (left) and
    the string diagram it denotes (right), and the labels beside each boundary name its typed ports.
    Both legs, order (blue) and item
    (red), carry the default count of one, so no constraint system is drawn. $g_a$ is shown in its $b$-redo context, the
    $c$-context is analogous, and the entry and exit contexts pair $a$ with the open ends $\bot,\top$.}
  \label{fig:re-loop-cospans}
\end{figure*}

The sequence and parallel operators of Definition~\ref{def:process-tree} have an exact
characterisation in terms of series-parallel partial orders, developed in
Appendix~\ref{subsec:pt-sp-appendix} (Theorem~\ref{thm:pt-as-sp-language},
Corollary~\ref{cor:pt-sp-iff}).
